\section{System Design}
\label{sec:design}

CPSForge is a closed-loop LLM red/blue team framework for CPS security analysis. It comprises five interacting layers---physical execution, observation, red team (attack generation), safety enforcement, and blue team (defense)---connected to a real PLC. The same LLM serves both roles with different prompts. The framework is configuration-driven: any PLC-connected plant can be analyzed by providing a scene YAML. Figure~\ref{fig:architecture} illustrates the pipeline.

\begin{figure*}[t]
    \centering
    \begin{tikzpicture}[
        >=Stealth,
        node distance=0.6cm and 0.5cm,
        every node/.style={font=\small},
        layer/.style={draw, rounded corners=3pt, minimum width=2.2cm, minimum height=0.7cm, align=center, thick},
        plcbox/.style={layer, fill=orange!15, minimum width=3cm},
        atkbox/.style={layer, fill=red!12},
        shieldbox/.style={layer, fill=yellow!20},
        defbox/.style={layer, fill=blue!12},
        logbox/.style={layer, fill=gray!12},
        llmbox/.style={layer, fill=purple!12},
        groupbox/.style={draw, dashed, rounded corners=5pt, inner sep=6pt, thick},
        arr/.style={->, thick, >=Stealth},
        darr/.style={<->, thick, >=Stealth},
    ]

    %% ===== Physical Execution Layer =====
    \node[plcbox] (plc) {Siemens S7\\PLC};
    \node[plcbox, right=1.0cm of plc] (fio) {Factory I/O\\(or any plant)};
    \draw[darr] (plc) -- (fio);

    %% ===== Observation Layer =====
    \node[logbox, below=0.8cm of $(plc)!0.5!(fio)$, minimum width=5.0cm] (mw) {Observation Layer\\(python-snap7, 500\,ms polling, phase inference)};
    \draw[arr] (plc) -- (mw);

    %% ===== Context Builder =====
    \node[logbox, below=0.6cm of mw, minimum width=5.0cm] (ctx) {Context Builder\\(minimal $|$ partial $|$ full)};
    \draw[arr] (mw) -- (ctx);

    %% ===== Red Team (left) =====
    \node[atkbox, below left=1.0cm and 1.5cm of ctx, minimum width=2.5cm] (llm_atk) {\textbf{LLM Red Team}\\(Online MITM)};
    \node[atkbox, below=0.25cm of llm_atk] (static_atk) {Static LLM};
    \node[atkbox, below=0.25cm of static_atk] (random_atk) {Random Baseline};

    %% Context to red team
    \draw[arr, red!60] (ctx.south) -- ++(0,-0.3) -| (llm_atk.north);

    %% ===== Defense Chain (center-right) — sequential =====
    \node[shieldbox, below=3.8cm of ctx, minimum width=3.0cm, minimum height=0.7cm] (shield) {\textbf{Safety Shield}\\(range, whitelist, invariant)};

    \node[defbox, right=0.4cm of shield, minimum width=2.2cm] (phase_shield) {PhaseAware\\Shield};
    \node[defbox, right=0.4cm of phase_shield, minimum width=2.2cm] (intent) {Intent\\Checker};
    \node[llmbox, right=0.4cm of intent, minimum width=2.2cm] (llm_def) {\textbf{LLM Blue}\\
    \textbf{Team}};

    %% Red team proposals → defense chain (sequential)
    \draw[arr, red!60] (llm_atk.south) -- ++(0,-0.2) -| (shield.north);
    \draw[arr, red!40] (static_atk.south) -- ++(0,-0.1) -| ([xshift=-0.3cm]shield.north);
    \draw[arr, red!40] (random_atk.south) -- ++(0,-0.05) -| ([xshift=-0.6cm]shield.north);

    %% Sequential defense chain arrows
    \draw[arr] (shield) -- (phase_shield);
    \draw[arr] (phase_shield) -- (intent);
    \draw[arr] (intent) -- (llm_def);

    %% Context to blue team LLM
    \draw[arr, blue!60] (ctx.south) -- ++(0,-0.3) -| (llm_def.north);

    %% ===== PLC Write or Block =====
    \node[plcbox, below=0.7cm of intent, minimum width=2.5cm] (write) {PLC Write\\(execute)};
    \node[logbox, below=0.7cm of llm_def, minimum width=2.0cm] (block) {Blocked\\(logged)};
    \draw[arr, green!60!black] (llm_def.south) -- ++(0,-0.15) -| (write.north) node[pos=0.25, right, font=\scriptsize]{allow};
    \draw[arr, red!60] (llm_def.south) -- (block) node[midway, right, font=\scriptsize]{block};

    %% Arrow back to PLC
    \draw[arr, thick, orange!60] (write.west) -- ++(-1.5,0) |- (plc.south);

    %% ===== Artifacts =====
    \node[logbox, below=0.5cm of write, minimum width=4.5cm] (artifacts) {Run Artifacts (Parquet, JSON, PLC event capture)};
    \draw[arr, gray] (block.south) |- (artifacts.east);

    %% ===== Group boxes =====
    \begin{scope}[on background layer]
        \node[groupbox, fill=orange!5, fit=(plc)(fio), label={[font=\scriptsize\bfseries]above:Physical Execution Layer}] {};
        \node[groupbox, fill=red!5, fit=(llm_atk)(static_atk)(random_atk), label={[font=\scriptsize\bfseries]above:Red Team Variants}] {};
        \node[groupbox, fill=blue!5, fit=(shield)(phase_shield)(intent)(llm_def), label={[font=\scriptsize\bfseries]above:Blue Team Defense Chain (sequential)}] {};
    \end{scope}

    \end{tikzpicture}
    \caption{CPSForge architecture. The observation layer polls the real PLC and builds context for both red and blue team LLMs. Attack proposals flow through the sequential defense chain: Safety Shield $\to$ PhaseAwareShield $\to$ IntentChecker $\to$ LLM Blue Team. A write must pass all stages to reach the PLC. The same model and context pipeline serve both roles.}
    \label{fig:architecture}
\end{figure*}

\subsection{Physical Execution Layer}

A Siemens S7-1200 PLC executes control logic deployed from TIA Portal~v17. Factory~I/O serves as the process emulator. Communication uses python-snap7 over Ethernet at 500\,ms polling. Each scene is defined by a YAML profile containing tag metadata (addresses, types, ranges), process descriptions, attack surface definitions, safety rules, and reset procedures. Three scenes cover continuous and discrete processes:

\begin{itemize}[leftmargin=1.2em,nosep]
    \item \textbf{Level Control} (15~tags): continuous PID tank fill/drain.
    \item \textbf{Sorting by Weight} (35~tags): discrete 3-way weight classification.
    \item \textbf{Sorting by Height} (30~tags): discrete 2-way height sorting.
\end{itemize}

\subsection{Observation and Context Layer}
\label{sec:design:context}

The middleware reads PLC tags every 500\,ms, constructing structured snapshots (sensor values, actuator states, setpoints, alarms, derived features). A rule-based phase inference engine determines the current operational phase from snapshot patterns. A modular context builder assembles three tiers for the RQ1 ablation:

\begin{itemize}[leftmargin=1.2em,nosep]
    \item \textbf{Minimal:} current tag values only. No history, descriptions, or ranges.
    \item \textbf{Partial:} values + 5-step history + tag descriptions + short scene summary + derived features (rates of change, moving averages).
    \item \textbf{Full:} values + 10-step history + inferred phase + prior action effects + attack guidance + safety rules.
\end{itemize}

The same context builder serves both red and blue team LLMs, ensuring a fair information basis.

\subsection{LLM Attacker Layer}
\label{sec:design:attacker}

Three attacker variants enable controlled comparison (RQ3):

\paragraph{Random Baseline.} Bounded random perturbations over the scene's attack surface. Provides the lower bound.

\paragraph{Static LLM.} One-shot batch: receives scene context and a single snapshot, proposes attacks without process feedback.

\paragraph{Online MITM (C1).} The core contribution. Every $k{=}3$ polling cycles: (1)~observe current PLC state; (2)~inspect history; (3)~infer phase; (4)~assemble context-tier prompt; (5)~LLM decides attack or wait; (6)~if attack, emit structured JSON with target tag, value, duration, reasoning; (7)~pass through defense chain; (8)~observe result; (9)~continue iteratively.

All LLM outputs must parse into a strict \texttt{AttackDecision} JSON schema. Invalid outputs default to ``wait.'' Attack budget: max 10 per run. The LLM never receives raw PLC addresses---only symbolic tag names.

\subsection{Safety Shield}
\label{sec:design:shield}

The mandatory first gate before any write reaches the PLC. It evaluates sequential checks: (1)~whitelist (tag must be in declared attack surface); (2)~range (value within configured bounds); (3)~duration cap; (4)~cooldown window; (5)~mutual exclusion (conflicting actuators cannot be simultaneously active); (6)~invariant (hard physical-boundary rules). The shield runs in microseconds and is always active regardless of defense variant. It also generates a rollback plan (original values for automatic restoration after attack expiry).

\subsection{Defense Layer}
\label{sec:design:defense}

Seven defense configurations (RQ4) organized into three categories:

\paragraph{Baseline.} Threshold and invariant detectors (rule-based).

\paragraph{Novel Rule-Based Defenses.} \emph{PhaseAwareShield}: blocks writes inconsistent with the current operational phase (e.g., opening discharge during fill). Standard range checks cannot catch contextually wrong but individually legal values. \emph{IntentConsistencyChecker}: blocks writes opposing the controller's current trend (e.g., forcing fill valve open while PID is reducing it).

\paragraph{LLM Defender (C4).} Uses the \emph{same model, context builder, and inference pipeline} as the attacker, with a defense-oriented prompt. For each proposed write that passes the safety shield, it receives: current process state, proposed write details, scene description, inferred phase. It outputs: \texttt{\{decision: allow|block, suspicion\_score: 0--1, reasoning\}}. On parse failure, defaults to ``allow'' to avoid blocking legitimate writes. This creates a symmetric LLM-vs-LLM evaluation.

The full defense chain applies sequentially: Safety Shield $\to$ PhaseAwareShield $\to$ IntentChecker $\to$ LLM Defender $\to$ Execute. Each stage can independently block a write.

\subsection{Fine-Tuning Pipeline}
\label{sec:design:finetune}

QLoRA fine-tuning (C3) trains Qwen2.5-3B-Instruct on 2{,}057 examples: 83~post-hoc verified successful attack traces (tag moved toward injected value), 608~defense ``block'' examples (known attack proposals), and 1{,}366~defense ``allow'' examples (normal operation with synthetic legitimate writes). Configuration: rank=16, $\alpha$=32, dropout=0.05, 3~epochs, lr=$2{\times}10^{-4}$, 4-bit NF4 quantization. Trainable parameters: 29.9M (0.96\% of 3.1B).
